% Compile twice: lualatex 30004898.tex
\documentclass[11pt]{article}
\usepackage[margin=25mm]{geometry}
\usepackage{fontspec,unicode-math}
\setmainfont{Latin Modern Roman}
\setmathfont{Latin Modern Math}
\usepackage{amsmath,amssymb,mathtools}
\usepackage{xurl,hyperref}
\hypersetup{colorlinks=true,urlcolor=blue,linkcolor=blue}
\setcounter{secnumdepth}{0}
\setlength{\parindent}{0pt}
\setlength{\parskip}{5pt}
\setlength{\emergencystretch}{3em}
\newcommand{\sref}[2]{\hyperlink{source:#1:#2}{[S#2]}}
\newcommand{\cataloguescope}[1]{\paragraph{Source question.}#1}
\newcommand{\floor}[1]{\left\lfloor#1\right\rfloor}
\begin{document}
% UnsolvedMath ID 30004898; source catalogue code OWR-8415356-002
\setcounter{section}{30004897}
\section{Complexity of Submodular $k$-Partition}\label{30004898}
{\small\sffamily UnsolvedMath ID 30004898 · Computer Science}
\subsection{Definitions and mathematical statement}

\paragraph{Shared notation.}Standard number systems are denoted by $\mathbb N,\mathbb Z,\mathbb Q,\mathbb R,\mathbb C$; index ranges, computational models and specialized notation follow the cited problem source and the definitions below.


A set function $f:2^V\to\mathbb R$ is submodular if $f(S)+f(T)\geq f(S\cup T)+f(S\cap T)$ for all $S,T\subseteq V$. A $k$-partition is a collection of pairwise disjoint nonempty sets $V_1,\ldots,V_k$ with union $V$. Its cost is $\sum_{i=1}^kf(V_i)$. The integer $k$ is fixed while $|V|$ grows.
\paragraph{Statement recorded in the supplied source.}
For fixed $k$, what is the computational complexity of minimizing the sum of a submodular function over a partition into $k$ nonempty parts? \sref{30004898}{2}
\paragraph{Scope qualification.}
Complexity depends on the representation of $f$; a value-oracle model and an explicit finite representation are different problems. The supplied question does not impose monotonicity, symmetry, or a particular representation, so none is assumed here.
\paragraph{Source status.}
The supplied catalogue labels this record partially solved. Its source-status metadata was checked against the supplied snapshot on 2026-10-11. This import does not independently establish current open status or mathematical correctness. The catalogue literature assessment is dated 2026-08-22.
\subsection{Short English statement}
Determine the complexity of dividing a ground set into a fixed number of nonempty parts of minimum total submodular cost.
\subsection{Sources}
\begin{itemize}
\item[\hypertarget{source:30004898:1}{S1}] Ulam AI and the UnsolvedMath Contributors, UnsolvedMath: A Curated Collection of Open Mathematics Problems, record 30004898 (OWR-8415356-002). Catalogue curation and metadata: CC BY 4.0. The supplied source question is retained; definitions and the English summary are editorial additions.
\url{https://huggingface.co/datasets/ulamai/UnsolvedMath}.
\item[\hypertarget{source:30004898:2}{S2}] Combinatorial Optimization (2022), Oberwolfach Reports, DOI 10.4171/owr/2021/53, pp. 2893--2954
\url{https://doi.org/10.4171/owr/2021/53}.
\item[\hypertarget{source:30004898:3}{S3}] Combinatorial Optimization, Oberwolfach Reports 53/2021, primary report PDF
\url{https://ems.press/content/serial-article-files/46931}.
\textit{Underlying Oberwolfach Reports material retains its source-specific licence, including CC BY-SA 4.0 where applicable; see THIRD\_PARTY\_NOTICES.md.}
\end{itemize}
\paragraph{Preparation and review.}
Original question attributed to the cited source and UnsolvedMath contributors. Codex (GPT-6; exact serving revision not exposed) assisted with definitions, English summaries, source inspection and formatting. No new human mathematical authorship or independent proof verification is asserted. Canonical ID 30004898 and source code OWR-8415356-002 checked against the complete supplied registry and existing repository statements; known flagged repetitions and obvious identical conjectures excluded. Full mathematical-equivalence/literature review has not been performed. Statement preparation and final review followed contribution prompts C and D; automated checks verify structure and identity, not open status.
\label{end:30004898}
\end{document}
